\chapter{Implementierung}

\section{Entwicklungs- und Ausführungsumgebung}

Die experimentelle Pipeline wurde Python-basiert in einer GPU-gestützten
Kaggle-Notebook-Umgebung umgesetzt. Für Training und Inferenz der neuronalen
Modellkomponenten wird PyTorch verwendet, während die transformerbasierten Encoder und
Tokenizer über die Transformers-Bibliothek eingebunden werden. Ergänzend kommen NumPy,
pandas und PyArrow für die Datenverarbeitung sowie scikit-learn und SciPy für
Downstream-Modelle und statistische Auswertungen zum Einsatz. Das statische HTML- und
DOM-Parsing der vorgelagerten Data-Freeze-Pipeline basiert auf lxml.

Die rechenintensiven Verarbeitungsschritte werden CUDA-beschleunigt ausgeführt. Die finale
Integrationspipeline prüft hierzu beim Start die Verfügbarkeit einer GPU und verwendet eine
NVIDIA-Tesla-T4-Umgebung als Zielkonfiguration. Für Trainings- und Inferenzoperationen ist
darüber hinaus Mixed Precision aktiviert; die entsprechenden Berechnungen werden über
\texttt{torch.autocast} mit \texttt{float16} ausgeführt. Speicherintensive
Verarbeitungsblöcke enthalten zusätzliche CUDA-Speicherbereinigungen. Für mehrere Trainings-
und Scoring-Schritte sind zudem abgestufte Batchgrößen implementiert, sodass ein
Verarbeitungsschritt bei einem \texttt{OutOfMemoryError} mit einer kleineren
Batchkonfiguration erneut ausgeführt werden kann. Die jeweils verwendeten Batchgrößen werden
bei den zugehörigen Modellkomponenten konkretisiert.

Für die reproduzierbare Durchführung der wiederholten Experimente werden die verwendeten
Zufallsquellen zentral initialisiert. Der Gesamtlauf verwendet den Master-Seed 20260817; die
einzelnen Versuchsblöcke greifen ergänzend auf festgelegte Modell- und Ranking-Seeds zurück.
Die Initialisierung umfasst Python, NumPy und PyTorch einschließlich der CUDA-Zufallsquellen.

Die während der Ausführung erzeugten Modellzustände, vorberechneten Repräsentationen,
Vorhersagescores sowie Ergebnis- und Auditinformationen werden getrennt als Artefakte
gespeichert. Die darauf aufbauenden Completion-, Freeze- und Reproduzierbarkeitsmechanismen
werden in Abschnitt 5.6 beschrieben.

\section{Implementierung der Datenverarbeitung}

Die Datenverarbeitung bildet den gemeinsamen Ausgangspunkt sämtlicher nachfolgenden
Versuchsblöcke. Der PhreshPhish-Datensatz wird zunächst in einem festen Data Freeze geprüft,
aufbereitet und in die in Abschnitt 4.2.1 definierten Rollen SUP, DEV, CAL und SSL überführt.
Der offizielle Testbestand bleibt davon getrennt und wird erst in den späteren
Evaluationsschritten als FINAL verwendet. Während des Data Freeze werden unter anderem die
erwarteten Rollenstärken, doppelte SHA-256-Identifikatoren und exakte Überschneidungen
zwischen Trainings- und Testbestand programmatisch geprüft.

Die Rollenbildung erfolgt deterministisch über die SHA-256-Identifikation der einzelnen
Webseiten. Hierzu wird der jeweilige Identifikator gemeinsam mit einem festen,
rollenspezifischen Zusatzstring erneut gehasht. Im Code wird dieser Zusatz als \texttt{salt}
bezeichnet. Dadurch entstehen für die unterschiedlichen Auswahlstufen getrennte, bei erneuter
Ausführung jedoch identische Rangfolgen. Listing~\ref{lst:stable_rank} im Anhang zeigt die hierfür
verwendete Funktion.



SUP und DEV werden auf dieser Grundlage jeweils klassenweise gebildet. CAL wird
anschließend ausschließlich aus noch nicht zugeordneten benignen Instanzen bestimmt, bevor
der SSL-Pool aus dem verbleibenden Bestand abgeleitet wird. Die Implementierung prüft
abschließend die vorgesehenen Größen von 4.000 SUP-, 20.000 DEV-, 50.000 CAL- und
200.000 SSL-Instanzen. Die deterministische Materialisierung setzt damit die in
Abschnitt 4.2.1 konzipierte Trennung der Datenrollen technisch um. Parallel zur Rollenbildung
werden die in Abschnitt 2.1.2 unterschiedenen Informationsperspektiven für die nachfolgenden
Modellzweige aufbereitet.

Für den textuellen HTML-Zweig werden zunächst Kommentare sowie \texttt{script}-,
\texttt{style}-, \texttt{noscript}- und \texttt{template}-Bereiche entfernt.
Anschließend werden die verbleibenden Tags entfernt, HTML-Zeichenreferenzen aufgelöst und
Whitespace normalisiert. Aus den normalisierten HTML-Tags wird zusätzlich eine strukturelle
Sequenz gebildet. Diese dient unter anderem zur Ableitung des \nolinkurl{template_hash}, der
später für die Template-basierten OOD-Szenarien benötigt wird.
Listing~\ref{lst:process_row} im Anhang fasst die für Text-, Template- und DOM-Ableitung maßgeblichen
Verarbeitungsschritte zusammen.



Der Ausschnitt zeigt nur die für die beschriebenen Ableitungen relevanten Teile der Funktion;
die vollständige Implementierung verbleibt in den digitalen Artefakten. Die DOM-Struktur wird
aus dem gespeicherten HTML statisch rekonstruiert. Für jeden berücksichtigten Knoten werden
Tagname, Elternindex, Tiefe, Attributanzahl und Anzahl direkter Kindknoten gespeichert. Damit
bleibt die in Abschnitt 2.2.2 beschriebene graphbasierte Sicht auf die
Eltern-Kind-Beziehungen des HTML-Dokuments für den späteren DOM-Zweig erhalten. Im Data
Freeze werden maximal 4.096 DOM-Knoten pro Webseite materialisiert; die engere
Eingabegrenze des Deep-Modells wird erst in Abschnitt 5.4 angewendet.

Eine zentrale technische Trennung betrifft den SSL-Pool. Entsprechend dem in Abschnitt 2.3
beschriebenen Prinzip des selbstüberwachten Lernens werden dessen physische Parquet-Dateien
ohne Downstream-Label gespeichert. Die Spalte \texttt{label} wird bereits während der
Rollenmaterialisierung explizit entfernt, wie Listing~\ref{lst:remove_label} im Anhang zeigt.



Die finale Pipeline kontrolliert zusätzlich über das Parquet-Schema, dass die physischen
SSL-Dateien weiterhin keine Labelspalte enthalten. Die Klasseninformationen verbleiben
getrennt in einem privaten Rollenmanifest und werden erst für die späteren überwachten
Labelbudgetexperimente über die SHA-256-Identifikation wieder zugeordnet. Die TAPT-Verfahren
arbeiten dagegen ausschließlich auf den physisch label-freien Webseiteninformationen. Damit
wird die experimentelle Trennung zwischen verfügbarem ungelabeltem Trainingsmaterial und
verfügbarer überwachter Trainingsinformation nicht nur konzeptionell, sondern auch in der
Datenpipeline umgesetzt.

\section{Implementierung der Self-Supervised Repräsentationen}

Die in Abschnitt 2.3 eingeführte selbstüberwachte Repräsentationsanpassung wird getrennt für
die URL- und die HTML-Text-Modalität umgesetzt. Als Ausgangspunkt dienen mit
\texttt{bert-base-uncased} beziehungsweise \texttt{roberta-base} bereits vortrainierte
Transformerencoder. Beide Modelle werden mittels Masked Language Modeling auf den
200.000 physisch label-freien Instanzen des SSL-Pools weitertrainiert. Damit wird das in
Abschnitt 2.3.3 beschriebene Task-Adaptive Pretraining auf dem aufgabenspezifischen
Webseiten-Trainingspool umgesetzt, ohne hierfür das Downstream-Label Phishing
beziehungsweise Benign zu verwenden.

Für die Modalitäten gelten unterschiedliche Kontextgrenzen: Der URL-Zweig verarbeitet
maximal 128 BERT-Token einschließlich Spezialtoken. Längere Eingaben werden abgeschnitten.
Eine vollständige Abdeckung des URL-Bestands durch diese Grenze ist damit nicht belegt.

Für den HTML-Text wird dagegen ein festes Kontextbudget von 256 Tokens verwendet. Die
HTML-Dokumente besitzen gegenüber URLs wesentlich umfangreichere und stärker variierende
Sequenzen. Die Begrenzung dient daher insbesondere dazu, Speicherbedarf und Rechenaufwand
des transformerbasierten Trainings kontrollierbar zu halten. Sie ist nicht als Annahme zu
verstehen, dass innerhalb von 256 Tokens der vollständige HTML-Inhalt einer Webseite erfasst
wird. Von dieser Transformergrenze ist die vorgelagerte Datenaufbereitung zu unterscheiden:
Im Data Freeze wird der bereinigte HTML-Text zunächst mit bis zu 50.000 Zeichen
materialisiert; die Begrenzung auf 256 Tokens erfolgt erst bei der Eingabe in den Textencoder.

Für beide Modalitäten wird das TAPT über eine Epoche durchgeführt. Da die Encoder dabei
nicht neu initialisiert, sondern bereits vortrainierte BERT- beziehungsweise RoBERTa-Modelle
zusätzlich an den Webseiten-Trainingspool angepasst werden, bildet eine Epoche ein bewusst
begrenztes Adaptationsbudget für den zusätzlichen SSL-Schritt. Beide Zweige durchlaufen damit
den jeweils 200.000 Instanzen umfassenden SSL-Bestand einmal vollständig. Die Festlegung ist
als einheitliche experimentelle Vorgabe und nicht als Aussage über ein global optimales
Trainingsbudget zu verstehen.

Für das URL-TAPT werden 15\,\% MLM-Maskierung, AdamW, eine Lernrate von
$5\cdot 10^{-5}$, Weight Decay 0,01 und 6\,\% Warmup verwendet.
Listing~\ref{lst:url_tapt} im Anhang zeigt die dafür maßgebliche Trainingskonfiguration.



Der TAPT-Pfad liest ausschließlich URLs aus dem physischen SSL-Pool; als Batch-Fallback
sind 16 und 8 hinterlegt. Für den HTML-Text-Zweig wird derselbe MLM-Mechanismus mit
15\,\% Maskierung, $5\cdot 10^{-5}$ Lernrate, Weight Decay 0,01 und 6\,\% Warmup
verwendet. Aufgrund der längeren Textsequenzen sind dort Batch-Fallbacks von 16, 8 und 4
vorgesehen. Der vollständige Trainingscode beider Zweige ist in den digitalen Artefakten
enthalten.

Die getrennte Anpassung beider Modalitäten bildet die technische Grundlage des in
Abschnitt 4.3.1 beschriebenen $2\times2$-Versuchsdesigns. Hierzu werden jeweils Basis- und
TAPT-Encoder so kombiniert, dass der Anpassungseffekt beider Modalitäten getrennt variiert
werden kann. R0 verwendet für URL und HTML-Text die unveränderten Ausgangsmodelle, RU
ausschließlich den angepassten URL-Encoder, RT ausschließlich den angepassten Textencoder
und RUT die TAPT-Varianten beider Zweige. Die verwendeten Encoderquellen werden für jede
Variante in den Audit-Artefakten gespeichert.

Damit bleibt die nachgelagerte Repräsentationsbildung zwischen den vier Bedingungen
strukturell identisch; variiert wird ausschließlich die Herkunft des URL- beziehungsweise
Textencoders. Auf diese Weise kann der zusätzliche Nutzen der task-adaptiven Anpassung
zunächst unabhängig von den späteren Architekturkomponenten wie DOM-Zweig und
multimodaler Fusion untersucht werden.

Für die kontrollierten Downstream-Probes werden aus den Hidden States der jeweiligen
Transformer feste Repräsentationsvektoren erzeugt. Hierzu wird ein
attention-mask-basiertes Mean Pooling verwendet, sodass Padding-Token nicht in die
Mittelwertbildung eingehen. Die Summe der gültigen Hidden States wird dabei durch die Anzahl
der über die Attention Mask markierten Token geteilt.

Die URL- und Textrepräsentationen werden getrennt extrahiert und als
\texttt{float32}-Embeddings zwischengespeichert. Die beiden 768-dimensionalen Vektoren werden
anschließend zu der in Abschnitt 4.3.1 definierten 1.536-dimensionalen Gesamtrepräsentation
konkateniert. Im Probe-Pfad erfolgt danach keine gemeinsame Weiteroptimierung der
Transformerencoder mit dem Klassifikator. Auf den festen Repräsentationen werden stattdessen
eine logistische Regression als Linear Probe sowie ein kleines MLP trainiert. Damit wird der
Repräsentationseffekt von Veränderungen der späteren End-to-End-Architektur getrennt.

Die Probe-Experimente werden für die sieben in Abschnitt 4.3.2 definierten Labelbudgets von
2.000, 5.000, 10.000, 20.000, 50.000, 100.000 und 200.000 gelabelten Instanzen
ausgeführt. Innerhalb einer Wiederholung wird hierfür eine deterministische klassenweise
Rangfolge verwendet, aus der verschachtelte Teilmengen entstehen. Ein kleineres Budget ist
damit Bestandteil der jeweils größeren Budgetstufen derselben Wiederholung. Für R0, RU, RT
und RUT werden innerhalb eines gepaarten Vergleichs dieselben Instanzen verwendet, sodass
sich die Bedingungen hinsichtlich der verfügbaren Labels nicht unterscheiden. Die Lernkurven
werden über fünf gepaarte Wiederholungen erzeugt. Für das zentrale Budget von 20.000 Labels
wird zusätzlich der in Abschnitt 4.3.2 beschriebene N10-Lauf durchgeführt; die konkrete
Seedsteuerung und Reproduzierbarkeitslogik wird in Abschnitt 5.6 zusammengeführt.

Ergänzend wurde in einem vorgelagerten Entwicklungsversuch ein kontrastiver SSL-Pfad
implementiert. R2C und R3C kombinieren modalitätsinterne Ansichten mit
URL--HTML-Positivpaaren; CSAME verwendet ausschließlich Same-Modality-Paare. Für den
historischen Vergleich wurde eine Temperatur von 0,10 verwendet. Um eng verwandte
Webseiten nicht als reguläre Negativbeispiele zu behandeln, werden Kandidaten mit identischer
nichtleerer Domain oder identischem nichtleerem Template-Hash aus der jeweiligen
Negativmenge ausgeschlossen.

Der kontrastive Pfad stammt aus einem eigenständigen früheren Versuchsblock und wird nicht
mit der finalen R0/RU/RT/RUT-Ablation vermischt. Seine Ergebnisse werden in Kapitel 6
deskriptiv als explorativer Objective-Vergleich dokumentiert; die finalen TAPT- und
Systementscheidungen werden daraus nicht abgeleitet.

\paragraph{Labelbudgetvergleich und klassische Referenz.}
Für den statistischen Labelbudgetvergleich werden die bereits extrahierten R0- und
RUT-Repräsentationen wiederverwendet; ein erneutes Transformertraining findet in diesem
Versuchsblock nicht statt. Die 1.536-dimensionalen URL+Text-Repräsentationen werden an eine
logistische Regression mit $C=1$, \texttt{solver=liblinear},
\nolinkurl{class_weight=balanced} und \nolinkurl{max_iter=2000} übergeben. RUT wird über das
vollständige in Abschnitt 4.3.2 festgelegte Budgetraster ausgewertet, während R0@200k als
interne Referenz mit identischem Repräsentationsformat und derselben Linear Probe dient.

Die ergänzende klassische Vergleichspipeline verwendet TF-IDF-Merkmale aus URL und Text.
Der URL-Pfad basiert auf Character-N-Grammen der Längen drei bis fünf mit
\nolinkurl{min_df=2}, sublinearer Termfrequenz, L2-Normalisierung und maximal 25.000
Merkmalen. Für den Text werden Word-N-Gramme der Längen eins bis zwei mit
\nolinkurl{min_df=3}, \nolinkurl{max_df=0.995}, sublinearer Termfrequenz,
Unicode-Akzentstripping, L2-Normalisierung und maximal 50.000 Merkmalen erzeugt. Die
TF-IDF-Transformation wird auf den 200.000 hierfür vorgesehenen Webseiten ohne Verwendung
der Klassenlabels angepasst; der resultierende Merkmalsraum umfasst 75.000 Dimensionen.

Die XGBoost-Referenz verwendet 400 Bäume mit maximaler Tiefe sechs, Lernrate 0,08,
\nolinkurl{min_child_weight=2}, \texttt{subsample=0.9},
\nolinkurl{colsample_bytree=0.65}, $\lambda=2$, \nolinkurl{max_bin=128} und
\nolinkurl{tree_method=hist}. Zehn Modelle mit getrennten Seeds werden im selben Runtime-Stand
neu trainiert und erhalten jeweils das vollständige 200.000er-Labelbudget. Die
Schwellenbestimmung erfolgt analog zum RUT/R0-Pfad auf CAL. Da sich gegenüber RUT+Linear
sowohl Repräsentationsverfahren als auch Klassifikator ändern, bleibt dieser Vergleich ein
Systembenchmark und wird nicht als isolierter TAPT-Effekt ausgewertet.

\section{Implementierung des multimodalen Deep-Systems}

Für das prototypische Gesamtsystem werden die in Abschnitt 5.3 ausgewählten Encoder in
einen gemeinsamen trainierbaren Modellpfad überführt. DUAL verarbeitet URL und HTML-Text,
TRI ergänzt die strukturelle DOM-Repräsentation.

Im Deep-System werden die obersten vier Transformerlayer für das überwachte Training
freigegeben. Listing~\ref{lst:freeze_last} im Anhang zeigt die entsprechende Parameterselektion.



Alle übrigen Transformerparameter bleiben eingefroren; ein vorhandener Pooler wird ebenfalls
trainierbar gehalten. Die über \texttt{rep} gewählten Basis- oder TAPT-Encoder werden nach
dem Masked Mean Pooling jeweils auf 256 Dimensionen projiziert und L2-normalisiert. Der
DOM-Zweig verarbeitet die im Data Freeze gespeicherten Knotenmerkmale und
Eltern-Kind-Beziehungen; pro Webseite werden maximal 512 Knoten in das Deep-System
übernommen. Listing~\ref{lst:dom_encoder} im Anhang zeigt den Kern der GCN-basierten Knotenkodierung.



Tag, Tiefe, Attribut- und Kindanzahl werden mit 96, 16, 8 und 8 Dimensionen eingebettet
und zu 128-dimensionalen Knotenvektoren kombiniert. Nach drei GCN-Schichten werden die
Knoten lernbar gewichtet aggregiert; der resultierende DOM-Vektor wird anschließend ebenfalls
auf 256 Dimensionen projiziert und L2-normalisiert.

Vor dem überwachten TRI-Training wird der DOM-Encoder über die Rekonstruktion maskierter
DOM-Tags selbstüberwacht vorangepasst.



Maskiert werden 30\,\% der Tags; trainiert wird über eine Epoche mit Batchgröße 64 und
$2\cdot10^{-3}$ Lernrate. Der Verlust wird ausschließlich auf den tatsächlich maskierten
Knoten bestimmt. Der resultierende DOM-Encoder wird gespeichert und im TRI-System
gemeinsam mit den übrigen Komponenten weiteroptimiert. Das strukturelle Rekonstruktionsziel
bleibt dabei vom tokenbasierten MLM der Transformerzweige getrennt.

Die aktiven Modalitätsrepräsentationen werden in \texttt{DeepFusion} zusammengeführt.
Listing~\ref{lst:deep_fusion} im Anhang zeigt den für Gating, Fusion und Main Head relevanten
Architekturausschnitt.



Das in Abschnitt 4.5.1 konzipierte Gating wird damit unmittelbar auf den normalisierten
Modalitätsvektoren umgesetzt. Für jede Modalität berechnet ein zweischichtiges Netz der Form
$256 \rightarrow 64 \rightarrow 1$ zunächst einen skalaren Gate-Wert. Ein Softmax über die
vorhandenen Modalitäten transformiert diese Werte in relative Gewichte. Die gemeinsame
Fusionsrepräsentation ergibt sich anschließend als gewichtete Summe:

\[
z_{\mathrm{Fusion}}
=
\sum_{m=1}^{M}\alpha_m z_m,
\qquad
\sum_{m=1}^{M}\alpha_m=1.
\]

Die Gewichte werden damit für jede Eingabe neu aus deren Modalitätsrepräsentationen
bestimmt. Neben diesem regulären gating-basierten Forward-Pfad enthält die finale
Integrationsimplementierung eine ungewichtete Mittelwertfusion für diagnostische
Auswertungen. Von diesem Auswertungsmodus ist die nachfolgende N10-Komponentenablation zu
unterscheiden: Dort werden die Equal-Varianten als eigenständige Modelle ohne lernbares Gate
trainiert.

Die fusionierte Repräsentation ersetzt die Einzelmodalitäten nicht. Stattdessen erhält der Main
Head sowohl die ursprünglichen Modalitätsvektoren als auch $z_{\mathrm{Fusion}}$. Für DUAL
entsteht dadurch der 768-dimensionale Eingabevektor

\[
[z_{\mathrm{URL}};z_{\mathrm{TEXT}};z_{\mathrm{Fusion}}],
\]

während TRI mit

\[
[z_{\mathrm{URL}};z_{\mathrm{TEXT}};z_{\mathrm{DOM}};z_{\mathrm{Fusion}}]
\]

einen 1.024-dimensionalen Eingang verwendet. Der Klassifikationskopf reduziert diese
Darstellung zunächst auf 512 und anschließend auf 128 Hidden Units. Zwischen den linearen
Schichten werden GELU-Aktivierungen und Dropout mit 0,2 beziehungsweise 0,1 eingesetzt,
bevor ein einzelner Hauptlogit erzeugt wird.

Zusätzlich besitzt jede vorhandene Modalität einen eigenen linearen Auxiliary Head. Diese
Köpfe arbeiten unmittelbar auf den finalen normalisierten 256-dimensionalen Vektoren
$z_{\mathrm{URL}}$, $z_{\mathrm{TEXT}}$ und gegebenenfalls $z_{\mathrm{DOM}}$. Ihre
Ausgaben werden erst nach der Berechnung des Main Heads erzeugt und gehen weder in die
Bestimmung der Gate-Gewichte noch in die Fusion oder als Eingabemerkmale in den Main Head
ein. Für den URL-Zweig erzeugt der lineare Auxiliary Head den skalaren Logit

\[
f_{\mathrm{URL}}
=
w_{\mathrm{URL}}^{\top}z_{\mathrm{URL}} + b_{\mathrm{URL}}.
\]

Höhere Werte entsprechen einer stärkeren Zuordnung zur Phishing-Klasse, niedrigere Werte
einer stärkeren Zuordnung zur benignen Klasse. $f_{\mathrm{URL}}$ wird nicht als kalibrierte
Wahrscheinlichkeit interpretiert, sondern als URL-spezifischer Klassifikationsscore. Dieser
Score wird in Abschnitt 5.5 als Routinggröße der URL-first-Kaskade wiederverwendet. Die
Auxiliary-Ausgaben bleiben damit vom Main Head und von der Fusionsentscheidung strukturell
getrennt.

Die Auxiliary Heads wirken dennoch auf die Parameteroptimierung, da ihre
Klassifikationsverluste als zusätzliche Trainingssignale in die Gesamtzielfunktion eingehen.
Der Hauptverlust wird um 30\,\% des mittleren Auxiliary-Loss ergänzt. Die
Auxiliary-Verluste werden damit als zusätzliches Trainingssignal verwendet, ohne die Struktur
des Main Heads oder die Berechnung der Gate-Gewichte zu verändern.

Für die Optimierung werden Transformerencoder und die übrigen Systemkomponenten in
getrennten Parametergruppen geführt. Die freigegebenen URL- und Textencoder werden mit
einer Lernrate von $1\cdot10^{-5}$ weiteroptimiert. Für Projektions-, DOM-, Gate-,
Auxiliary- und Klassifikationskomponenten wird dagegen $2\cdot10^{-4}$ verwendet. Beide
Gruppen nutzen einen Weight Decay von 0,01; der lineare Warmup umfasst 5\,\% der
vorgesehenen Optimierungsschritte. Die beiden Parametergruppen werden gemeinsam mit AdamW
optimiert.

Das reguläre Deep-Training umfasst eine Epoche, Mixed Precision und Gradient Clipping mit
maximaler Norm 1,0; als Batch-Fallbacks sind für DUAL 16/8/4 und für TRI 8/4/2 vorgesehen.
Die ergänzende Komponentenablation verwendet davon getrennt das nachfolgende feste
Kontrollprotokoll.

\subsection{Sequenzielle Komponentenablation}

Für die isolierte FF4-Prüfung wird die Systemfamilie in vier verschachtelte Varianten zerlegt:
P0\_URL enthält nur den task-adaptierten URL-Zweig, P1\_DUAL\_EQUAL ergänzt Text mit
gleichgewichteter Fusion, P2\_TRI\_EQUAL zusätzlich den DOM und P3\_TRI\_GATED dieselben
drei Modalitäten mit lernbarem Gating.

Die initialen RUT-URL- und Textcheckpoints sowie der DOM-SSL-Encoder sind über die
Replikationen identisch; die im Deep-Training freigegebenen Parameter werden weiterhin
überwacht optimiert. Pro Replikation
erhalten alle vier Varianten dieselbe balancierte Stichprobe aus 20.000 Labels; zehn Modell-
und zehn korrespondierende Ranking-Seeds bilden die gepaarten N10-Läufe. Jede Variante wird
eine Epoche mit festem Batch acht trainiert, sodass kein zusätzlicher Unterschied durch
OOM-Fallbacks entsteht. Die Schwelle für 0,5\,\% Ziel-FPR wird ausschließlich auf den
benignen Beispielen von ENG\_TUNE bestimmt, die primäre TPR auf dem disjunkten
ENG\_META. CAL und FINAL werden nicht aufgerufen. Protokolliert werden die drei
angrenzenden TPR-Kontraste Text, DOM und Gating; Gate-Gewichte dienen nur der ergänzenden
Diagnose.

\section{Implementierung der Kaskade und Messinstrumentierung}

Die URL-first-Kaskade nutzt den in Abschnitt 5.4 definierten URL-Auxiliary-Score
$f_{\mathrm{URL}}$ zur Routingentscheidung. Der Score ist ein unkalibrierter Logit des
URL-Auxiliary-Heads und dient innerhalb der Kaskade ausschließlich zur Rang- und
Schwellenentscheidung. Eine direkte URL-Entscheidung wird nur in den beiden äußeren
Bereichen der auf Development beobachteten Scoreverteilung zugelassen. Instanzen im
dazwischenliegenden Routingintervall werden an den vollständigen multimodalen Pfad
eskaliert.

Die beiden Routinggrenzen werden ausschließlich auf Development-Daten festgelegt.
Listing~\ref{lst:cascade_thresholds} im Anhang zeigt deren Berechnung.



Die obere Grenze $t_{\mathrm{high}}$ wird anhand der benignen URL-Scores so bestimmt, dass
sie einem FPR-Ziel von 0,1\,\% auf Development entspricht. Scores oberhalb dieser Grenze
werden direkt der Phishing-Klasse zugeordnet. Die untere Grenze $t_{\mathrm{low}}$ wird aus
dem unteren Prozent der positiven URL-Scoreverteilung abgeleitet. Scores unterhalb dieser
Grenze werden direkt als benign eingeordnet. Beide Grenzen sichern damit unterschiedliche
direkte Entscheidungen ab: $t_{\mathrm{high}}$ kontrolliert den oberen, phishingnahen
Scorebereich anhand der benignen Instanzen, während $t_{\mathrm{low}}$ den unteren,
benignnahen Bereich anhand der positiven Scoreverteilung abgrenzt.

Für eine Instanz $x$ ergibt sich die Routingentscheidung damit als

\[
\operatorname{route}(x)=
\begin{cases}
\text{Benign}, &
f_{\mathrm{URL}}(x)\leq t_{\mathrm{low}},\\[3pt]
\text{TRI-Full-Pfad}, &
t_{\mathrm{low}} < f_{\mathrm{URL}}(x) < t_{\mathrm{high}},\\[3pt]
\text{Phishing}, &
f_{\mathrm{URL}}(x)\geq t_{\mathrm{high}}.
\end{cases}
\]

Der Bereich zwischen beiden Grenzen wird als Unsicherheitsbereich der Kaskade bezeichnet.
Dabei handelt es sich nicht um ein probabilistisch kalibriertes Unsicherheitsmaß, sondern um
ein empirisch auf Development bestimmtes Scoreintervall, innerhalb dessen keine direkte
URL-basierte Entscheidung getroffen wird. Die Routinggrenzen sind zudem vom später auf CAL
bestimmten 0,5-\%-Betriebspunkt des finalen Klassifikators zu unterscheiden.

Listing~\ref{lst:cascade_route} im Anhang zeigt die entsprechende Routinglogik.



Die Kaskade setzt damit nicht voraus, dass der URL-Zweig über den gesamten Datenbestand
dieselbe Klassifikationsleistung wie das multimodale Vollsystem erreicht. Der URL-Score wird
nur für die auf Development abgegrenzten äußeren Scorebereiche als direkte
Entscheidungsgrundlage verwendet. Für den verbleibenden Zwischenbereich bleibt die
vollständige multimodale Verarbeitung erhalten. Dadurch können eindeutige URL-Fälle früh
abgeschlossen werden, ohne HTML-Text, DOM und Fusion für alle Instanzen auszuführen.

Für eskalierte Instanzen wird die bereits berechnete URL-Repräsentation nicht erneut erzeugt.
Stattdessen wird $z_{\mathrm{URL}}$ für die betroffenen Zeilen wiederverwendet und um den
HTML-Text-Zweig sowie in der TRI-Variante den DOM-Zweig ergänzt. Die entstehenden
Modalitätsvektoren durchlaufen anschließend dieselbe gating-basierte Fusion und denselben Main
Head wie im vollständigen Deep-System. Die Kaskade verändert damit weder die zugrunde
liegenden Repräsentationen noch die Fusionslogik, sondern ausschließlich den für eine Instanz
tatsächlich ausgeführten Berechnungspfad.

Bei der Reproduktion sind Datentyp und Grenzwertbehandlung Teil des Protokolls:
die gespeicherten Kaskadenscores beruhen auf Float32-Logits und der damaligen
NumPy-Skalarumwandlung. Eine nachträgliche Umstellung des Routingvergleichs auf Float64
ändert bei gebundenen Scores die Eskalationsmenge. Das digitale Prüfskript setzt die
Routinggrenzen deshalb explizit auf Float32, während der abschließende
CAL-Klassifikationsvergleich wie im Originalcode in Float64 erfolgt.

Die Kaskadenkonfiguration wird auf Development-Daten festgelegt; kontrolliert werden
TPR-Verlust, Eskalationsrate und Fehlalarmrate gegenüber dem Vollpfad. Erst diese
Konfiguration geht in den Architecture-Freeze ein, sodass CAL und FINAL nicht zur
nachträglichen Optimierung der Routinggrenzen verwendet werden.

Nach drei Warm-up-Durchläufen wird die Einzelinstanzlatenz für 500 Webseiten bei Batchgröße
eins gemessen; CUDA wird vor der Zeitmessung synchronisiert. Berichtet werden Median und
95. Perzentil. Der Durchsatz wird auf 5.000 Webseiten bei Batchgröße 32 über drei
Wiederholungen bestimmt; zusätzlich wird Peak-VRAM erfasst. Der Benchmark umfasst nur den
aufbereiteten In-Memory-Modellpfad ohne Netzwerkzugriffe, Browser-Rendering,
Datenträger-I/O und vorgelagerte HTML-Verarbeitung.

Die Kaskade wird zusätzlich mit \nolinkurl{URL_BASE} und \nolinkurl{URL_TAPT} als Frontend
ausgeführt; TRI-Fallback und Routinglogik bleiben unverändert. Auf Development-Daten werden
Routingparameter für erlaubte TPR-Abweichungen von primär 1,0 sowie 0,25, 0,5 und
2,0 Prozentpunkten bestimmt. Fünf Systemläufe werden über alle fünf finalen Bedingungen
bewertet. Zentrale Kennzahl ist die Eskalationsrate beziehungsweise der Anteil vermiedener
Full-Path-Ausführungen; die TPR-Vorgabe beschreibt die auf Development festgelegte Toleranz
und keine identische realisierte Test-TPR.

\section{Artefaktbereitstellung und Reproduzierbarkeit}

Seeds, deterministische Datenpartitionierungen sowie gespeicherte Checkpoints, Embeddings,
Scores und Auditdateien dokumentieren die einzelnen Verarbeitungsschritte. Bereits
abgeschlossene Schritte werden durch Completion-Marker gekennzeichnet, sodass ein
unterbrochener Lauf an definierten Stellen fortgesetzt werden kann. Die Datei
\nolinkurl{ARCHITECTURE_FREEZE.json} hält die auf Development getroffenen
Systementscheidungen fest, bevor CAL und FINAL verwendet werden. Beim Übergang in die
abschließende Evaluationsphase wird zusätzlich erneut die SHA-256-Disjunktheit der
Datenrollen geprüft.

Für die ergänzte Labelbudgetanalyse und die FF4-Komponentenablation werden Analyseplan,
Seeds, Testgruppen und Schwellenlogik vor der jeweiligen Ausführung gespeichert.
Anschließende Coverage-Audits prüfen, ob die erwarteten Ergebnisdateien für alle vorgesehenen
Varianten und Wiederholungen vorliegen. Bei fehlenden Ergebnissen wird die Auswertung
abgebrochen. Dadurch lassen sich die in Kapitel 6 berichteten Ergebnisse ihren jeweiligen
Konfigurationen und Versuchsständen eindeutig zuordnen.


